\subsection{Agreed departure rule: neighbor information only}
For a cell $i$ considering departure, write $W_i$ for its current weight,
$N_i$ for its adjacent cells, $d_i=|N_i|$, and
\[
 m_i=\operatorname{median}_{j\in N_i} W_j.
\]
The median excludes $i$; for an even number of neighbors it is the average
of the two middle weights. The center departs when both conditions hold:
\[
 \boxed{W_i<\min_{j\in N_i}W_j\quad\hbox{and}\quad
 W_i<\frac{d_i}{d_i+1}m_i.}
\]
Thus it must be strictly lighter than every neighbor. Equality does not trigger
departure. No tunable coefficient $\alpha$ remains. An empty neighbor set is
an invalid decision state and does not authorize departure.

The second condition is equivalent to
\[
 m_i-W_i>\frac{W_i}{d_i}.
\]
It compares the departing cell's deficit below the neighbor median with the
average added weight per neighbor if its weight were shared equally. Actual
Voronoi redistribution is not generally equal, so this motivates the threshold
without guaranteeing improved balance. The decision requires only the cell's
own weight and adjacent weights; no global ranking or global median is used.
Weights retain their density-model meaning. Candidate geometric changes are
not accepted or rejected by recounting pickup records.

\subsection{Committed departure and neighbor-to-neighbor migration}
The center notifies its neighbors and becomes inactive in the Voronoi diagram.
Its former territory redistributes among surviving sites. This is site deletion,
\emph{not} a requirement that its cell merge exactly with one particular neighbor.
The inactive center is a free-center token that can be passed between hosts;
it does not generate a cell while migrating.

The discussed routing rule starts at a heaviest former neighbor. Each host
compares its weight with those of its own adjacent cells and forwards the token
to a heaviest neighbor only if that neighbor is strictly heavier than itself.
Otherwise it is a local maximum, including a plateau maximum, and accepts the
free center for insertion. Every routing decision uses one-hop information.
For fixed weights and adjacency, strict increase prevents cycles and gives
termination on a finite graph. Changing weights during concurrent routing do
not have that guarantee. Tie-breaking between equally heavy eligible hosts,
concurrency, and a local synchronization or cycle-handling convention remain
implementation details to settle; a global freeze is not assumed.

Departure and insertion are committed. There is no tentative relocation,
objective-based veto, or rollback to the original location. One inactive center
means one fewer active Voronoi site temporarily; insertion restores the count.
This logical rule does not remove the need for reliable token delivery and
local coordination in a distributed implementation.

\subsection{Agreed insertion: analytical geometry, no optimization}
Let $H$ be the destination host's bounded convex Voronoi cell, with positive
area $A$. Use its uniform frozen density. Compute its area centroid and
area covariance:
\[
 \mu=\frac{1}{A}\int_H x\,dx,\qquad
 \Sigma=\frac{1}{A}\int_H (x-\mu)(x-\mu)^{\mathsf T}\,dx.
\]
Polygon area and moment formulas evaluate these quantities using only the
cell's boundary vertices, not individual pickup locations. Let $u$ be a unit
eigenvector of the $2\times2$ matrix $\Sigma$ corresponding to its largest
eigenvalue. This is the long-axis direction. If both eigenvalues are equal,
use the fixed coordinate $x$-axis; reversing $u$ only exchanges the two halves.

Find the unique offset $t$ such that the line perpendicular to $u$ divides
$H$ into equal areas:
\[
 H_- = H\cap\{x:u^{\mathsf T}x\leq t\},\qquad
 H_+ = H\cap\{x:u^{\mathsf T}x\geq t\},\qquad
 |H_-|=|H_+|=A/2.
\]
There is no search over cut orientation: covariance determines it. To compute
$t$ analytically, sort the vertex projections onto $u$. Between consecutive
distinct projections the cross-section length is linear in the offset, so
cumulative area is quadratic (possibly linear). Evaluate the area at the
interval endpoints, identify the interval containing $A/2$, and solve its
linear or quadratic equation for the root in that interval. This is finite
geometric calculation, not numerical objective optimization. Degenerate
projection intervals are skipped; floating-point tolerances must be recorded.

Place the existing and arriving centers at the geometric centroids
\[
 c_- = \frac{2}{A}\int_{H_-}x\,dx,\qquad
 c_+ = \frac{2}{A}\int_{H_+}x\,dx.
\]
Assign the two center identities deterministically to these positions. For a
full-dimensional convex host, both centroids lie inside the host and are
distinct, on opposite sides of the cut. Recompute the affected Voronoi cells
and adjacency. This completes insertion; there is no gradient search,
backtracking, multistart search, or optimization of center separation.

\paragraph{What the construction does and does not guarantee.}
The two geometric halves have equal modeled weight under the host's frozen
uniform density. Their centroid sites need not have the chosen cut as their
perpendicular bisector. In addition, boundaries with other cells change.
Therefore the resulting Voronoi cells are generally not those two halves and
need not have equal weights or optimal circularity. The construction supplies
a deterministic geometric placement, not a claim of optimal balance.
Ordinary local moves may improve the inserted cells afterward.

The user specified that any split-specific grade concern only the two split
cells, not a sum over neighboring grades. The final insertion rule performs
no grade optimization or grade-based acceptance at all; the earlier proposed
joint optimization has been withdrawn.

\subsection{Distributed scope, superseded proposals, and status}
Host geometry is available from its clipped cell and adjacent-center data;
departure weights and routing decisions are likewise local. Maintaining a
correct Voronoi diagram after deleting or inserting sites can create new
adjacencies. A distributed implementation must discover and exchange these
through local messages, possibly over several hops. Knowledge of the old
neighbor list alone must not be assumed sufficient for every resulting
neighbor's complete geometry. No global directory or global optimization is
part of the agreed decision rule.

This revision supersedes the earlier lightest-pair merge test, tunable
$\alpha d_i m_i$ threshold, search over all equal-area cuts, joint split-grade
optimization, and tentative rollback proposal. Those were discussion
alternatives, not implemented results. A spatial frozen-density integration
model was also discussed, but its adoption for redistribution has not been
settled. In particular, the ordinary per-cell formula $\rho_j A_j$ does not
by itself specify how to initialize densities after identities disappear or
are created. Weight/density transfer and refresh conventions must be made
explicit before implementation, without substituting post-move dot counts
for modeled candidate weights.

The intended new experiment returns to the saved end-of-round-59 state,
before badness-weighted selection began, preserving all later runs and files.
\textbf{This update changes the document only. No departure, migration,
insertion, restart, or simulation has been implemented or run.}
The geometric insertion and parameter-free departure condition are agreed;
distributed concurrency, tie conventions, and weight accounting still require
precise implementation definitions.
